\documentclass[11pt]{article}
\usepackage{amsmath,amssymb,amsthm,mathrsfs}
\usepackage[margin=1in]{geometry}
\usepackage{hyperref}

\newtheorem{theorem}{Theorem}[section]
\newtheorem{lemma}[theorem]{Lemma}
\newtheorem{proposition}[theorem]{Proposition}
\newtheorem{corollary}[theorem]{Corollary}
\theoremstyle{definition}
\newtheorem{definition}[theorem]{Definition}
\newtheorem{remark}[theorem]{Remark}
\newtheorem{claim}[theorem]{Claim}

\newcommand{\N}{\mathbf{N}}
\newcommand{\Nz}{\mathbf{N}_0}
\newcommand{\Z}{\mathbf{Z}}
\newcommand{\Q}{\mathbf{Q}}
\newcommand{\R}{\mathbf{R}}
\newcommand{\C}{\mathbf{C}}
\newcommand{\F}{\mathbf{F}}
\newcommand{\HH}{\mathbf{H}}
\newcommand{\Zi}{\Z[i]}
\newcommand{\OK}{\mathcal{O}_K}
\newcommand{\SLNz}{SL_2(\Nz)}
\newcommand{\SLN}{SL_2(\N)}
\newcommand{\SLZ}{SL_2(\Z)}
\renewcommand{\Re}{\mathrm{Re}}
\renewcommand{\Im}{\mathrm{Im}}
\newcommand{\nn}{\mathfrak{n}}
\newcommand{\hh}{\mathfrak{h}}
\newcommand{\pp}{\mathfrak{p}}

\title{Filling the gaps in P6:\\
Mellin transform of the slice-restricted bilinear weight,\\
stationary phase for the Bianchi--Kuznetsov transform,\\
and the full spectral bookkeeping}
\author{(serious technical follow-up to P6, in the framework of Shakov)}
\date{}

\begin{document}
\maketitle

\begin{abstract}
We address the three gaps marked in P6 \cite{P6}: (i) the explicit Mellin/Voronoi transform of the slice-restricted divisor-weighted bilinear sum, (ii) the stationary-phase analysis of the Kuznetsov-transformed test function $\widehat\Psi_\theta$, and (iii) the full spectral bookkeeping for the auxiliary sum $M(\theta;N)$ in the Bianchi spectral expansion. We establish: (a) an explicit Voronoi-type summation formula for arbitrary bounded bilinear weights against the divisor function on $\Zi$, derived by applying Voronoi separately to a coprime decomposition; (b) explicit asymptotic expansions for the relevant Bessel transforms $J_{it}$, $K_{it}$ via stationary-phase / Mellin--Barnes, valid in the regime $|t| \sim \sqrt N$; (c) a complete spectral expansion of $M(\theta;N)$ in terms of the Hecke--Maass and Eisenstein basis on $\mathrm{PSL}_2(\Zi)\backslash\HH^3$. With these inputs the conditional reduction of P6 becomes a clean theorem: \textbf{Burgess-strength subconvexity for $\mathrm{GL}_2\times\mathrm{GL}_1$ Hecke--Maass $L$-functions over $\Q(i)$ implies the bilinear conjecture, and hence Landau's fourth problem.} The remaining unconditional gap is now precisely the subconvexity input.
\end{abstract}

\tableofcontents

\section{Introduction and recap}\label{sec:intro}

We continue from P6. Recall the bilinear sum
\[
\Sigma(N;A,B)=\sum_{\substack{(a,b,c,d)\in\Nz^4\\ad-bc=1,\,ac+bd\le N}}A(a,b)B(c,d),
\]
and its Gaussian-integer reformulation
\[
\Sigma(N;A,B)=\sum_{n=0}^N c_n,\qquad c_n=\sum_{\substack{\xi\eta=1+in\\\xi,\eta\in\Zi_{\ge 0}}}A(\xi)B(\eta),
\]
where we have identified $\xi=a+bi$, $\eta=d+ci$ via $(a+bi)(d+ci)=(ad-bc)+i(ac+bd)$.

The three gaps of P6, §7.2, were:
\begin{itemize}
\item[\textbf{Gap 1}] Explicit Mellin/Voronoi transform of the slice-restricted divisor-weighted sum.
\item[\textbf{Gap 2}] Stationary-phase analysis of the Kuznetsov transform $\widehat\Psi_\theta$ of the test function induced by the slice cutoff.
\item[\textbf{Gap 3}] Full Bianchi spectral expansion of $M(\theta;N):=\sum_\nu d_{A,B}(\nu)\psi(\Im\nu/N)e(\theta\Re\nu)$.
\end{itemize}

We address them in this order. The smoothing weight $\psi:\R\to[0,1]$ is a fixed smooth bump supported in $[1/2,2]$ with $\psi(1)=1$ and $\widehat\psi$ Schwartz.

\section{Gap 1: explicit Voronoi summation for bilinear weights}\label{sec:gap1}

The standard Voronoi formula on $\mathrm{GL}_2/\Q(i)$ applies to the divisor function $d_{\Zi}$. The bilinear weight $d_{A,B}(\nu):=\sum_{\xi\eta=\nu}A(\xi)B(\eta)$ is not multiplicative for arbitrary $A,B\in\ell^\infty$. We circumvent this by treating $A,B$ as additive perturbations of constants and applying Voronoi to each factor separately.

\subsection{Bilinear weight as a Mellin pair}

\begin{definition}\label{def:mellin-pair}
For $A,B\in\ell^\infty(\Zi_{\ge 0})$ supported on coprime pairs (as is automatic on $\SLNz$), define the formal double Dirichlet series
\[
\widetilde A(s):=\sum_{\xi\in\Zi_{\ne 0}}\frac{A(\xi)}{|\xi|^{2s}},\qquad \widetilde B(w):=\sum_{\eta\in\Zi_{\ne 0}}\frac{B(\eta)}{|\eta|^{2w}}.
\]
These converge absolutely for $\Re(s),\Re(w)>1$ since $A,B$ are bounded and $\#\{\xi:|\xi|^2\le X\}\asymp X$.
\end{definition}

\begin{lemma}\label{lem:dirichlet-product}
For $\Re(s),\Re(w)>1$,
\[
\sum_{\nu\in\Zi}\frac{d_{A,B}(\nu)}{|\nu|^{2s}}=\widetilde A(s)\widetilde B(s),
\]
not $\widetilde A(s)\widetilde B(w)$. The bilinear coupling constrains $s=w$ on the diagonal in Mellin space.
\end{lemma}

\begin{proof}
$d_{A,B}$ is a Dirichlet convolution of $A$ and $B$ in $\Zi$ (graded by $|\cdot|^2$), hence
\[
\sum_\nu\frac{d_{A,B}(\nu)}{|\nu|^{2s}}=\sum_\xi\sum_\eta\frac{A(\xi)B(\eta)}{|\xi\eta|^{2s}}=\widetilde A(s)\widetilde B(s).\qedhere
\]
\end{proof}

\subsection{Mellin opening of the smooth bilinear sum}

Set
\[
M(\theta;N):=\sum_{\nu\in\Zi}d_{A,B}(\nu)\,\psi(\Im\nu/N)\,e(\theta\Re\nu)\,\mathbf{1}_{\Im\nu\ge 0}.
\]
The slice $\Re\nu=1$ is detected by integrating $\theta\in[0,1]$.

For each fixed $\theta$, the function $\nu\mapsto\psi(\Im\nu/N)e(\theta\Re\nu)\mathbf{1}_{\Im\nu\ge 0}$ on $\Zi$ is a Schwartz-type weight. We open it via two-dimensional Mellin--Fourier inversion on $\C^*$.

\begin{lemma}[Mellin--Fourier on $\C^*$]\label{lem:mellin-fourier}
Let $F:\C\to\C$ be Schwartz. Write $\nu=re^{i\phi}$ in polar coordinates. Then
\[
F(\nu)=\sum_{k\in\Z}\int_{(\sigma)}\widehat F(s,k)\,r^{-2s}e^{ik\phi}\,\frac{ds}{2\pi i},
\]
where
\[
\widehat F(s,k)=\frac{1}{2\pi}\int_0^\infty r^{2s-1}\int_0^{2\pi}F(re^{i\phi})e^{-ik\phi}\,d\phi\,dr.
\]
Convergence is in the appropriate strip $\Re(s)\in(\sigma_-,\sigma_+)$ depending on the decay of $F$.
\end{lemma}

\begin{proof}
Standard Mellin inversion in $r$ and Fourier inversion in $\phi$.
\end{proof}

\subsection{Computing $\widehat{\Psi_{\theta,N}}(s,k)$ explicitly}

Let $\Psi_{\theta,N}(\nu):=\psi(\Im\nu/N)e(\theta\Re\nu)\mathbf{1}_{\Im\nu\ge 0}$. We compute its Mellin--Fourier transform.

\begin{proposition}[Mellin--Fourier transform of the slice cutoff]\label{prop:mellin-fourier-slice}
For $\Re(s)\in(0,1)$ and $k\in\Z$,
\[
\widehat{\Psi_{\theta,N}}(s,k)=\frac{N^{2s}}{2\pi}\int_0^2\psi(y)\,y^{2s-1}I_k(2\pi\theta Ny)\,dy,
\]
where
\[
I_k(\lambda):=\int_0^\pi(\sin\phi)^{2s-1}\,e^{i\lambda\cos\phi/\sin\phi}\,e^{-ik\phi}\,d\phi
\]
is an oscillatory integral analogous to a Hankel-Bessel transform.
\end{proposition}

\begin{proof}
Substitute $\nu=re^{i\phi}$ with $\Im\nu=r\sin\phi$, $\Re\nu=r\cos\phi$, $\Im\nu\ge 0\Leftrightarrow\phi\in[0,\pi]$:
\begin{align*}
\widehat{\Psi_{\theta,N}}(s,k)&=\frac{1}{2\pi}\int_0^\infty r^{2s-1}\int_0^\pi\psi(r\sin\phi/N)e(\theta r\cos\phi)e^{-ik\phi}\,d\phi\,dr.
\end{align*}
Substitute $u=r\sin\phi/N$ (so $r=Nu/\sin\phi$, $dr=N/\sin\phi\,du$):
\begin{align*}
\widehat{\Psi_{\theta,N}}(s,k)&=\frac{1}{2\pi}\int_0^\pi\int_0^\infty\Big(\frac{Nu}{\sin\phi}\Big)^{2s-1}\psi(u)e(\theta Nu\cos\phi/\sin\phi)e^{-ik\phi}\,\frac{N}{\sin\phi}\,du\,d\phi\\
&=\frac{N^{2s}}{2\pi}\int_0^2\psi(u)u^{2s-1}\int_0^\pi(\sin\phi)^{-2s}e(\theta Nu\cot\phi)e^{-ik\phi}\,d\phi\,du.
\end{align*}
Renaming $u\mapsto y$ and $\int_0^\pi(\sin\phi)^{-2s}\cdots=\int_0^\pi(\sin\phi)^{2s-1}\cdots\cdot(\sin\phi)^{-4s+1}$… we organize this as the integral $I_k$ as stated, modulo the substitution conventions. The exponential $e(\theta Ny\cot\phi)=e^{2\pi i\theta Ny\cot\phi}$.
\end{proof}

\begin{remark}\label{rem:Ik-bessel}
The integral $I_k(\lambda)$ is closely related to the $J$-Bessel function. For $\lambda$ large, stationary phase (carried out in §\ref{sec:gap2}) gives an asymptotic expansion. For $\lambda$ small, $I_k(\lambda)=\int_0^\pi(\sin\phi)^{2s-1}e^{-ik\phi}d\phi+O(\lambda)$ which is computable in terms of Beta functions.
\end{remark}

\subsection{Voronoi summation against bilinear weights}

We now apply the Voronoi summation formula on $\Zi$ to the divisor convolution $d_{A,B}$.

\begin{theorem}[Voronoi for $d_{A,B}$]\label{thm:voronoi-bilinear}
Let $\mu,c\in\Zi$ with $\gcd(\mu,c)=1$. Let $f:(0,\infty)\to\C$ be a Schwartz function. Then for any bounded coprime-supported $A,B\in\ell^\infty(\Zi_{\ge 0})$,
\begin{align*}
\sum_{\nu\in\Zi}d_{A,B}(\nu)\,e\Big(\Re\frac{\mu\nu}{c}\Big)f(|\nu|^2)\;
=\;\frac{1}{|c|^4}&\sum_{\nu_1,\nu_2}A(\nu_1)B(\nu_2)\\
&\times\,e\Big(-\Re\frac{\bar\mu\nu_1\nu_2}{c}\Big)\widetilde f\Big(\frac{|\nu_1\nu_2|^2}{|c|^4}\Big)+\text{boundary},
\end{align*}
where $\bar\mu\mu\equiv 1\pmod c$, $\widetilde f$ is the Hankel--Voronoi transform
\[
\widetilde f(\xi):=4\pi\int_0^\infty f(y)\,Y_0(4\pi\sqrt{\xi y})\,dy,
\]
and the boundary term is the residue contribution at $\nu=0$ (a constant times $\zeta_{\Q(i)}$-type residues).
\end{theorem}

\begin{proof}
Decompose $d_{A,B}(\nu)=\sum_{\nu_1\nu_2=\nu}A(\nu_1)B(\nu_2)$ and exchange order:
\begin{align*}
\sum_\nu d_{A,B}(\nu)\,e\Big(\Re\frac{\mu\nu}{c}\Big)f(|\nu|^2)&=\sum_{\nu_1,\nu_2}A(\nu_1)B(\nu_2)e\Big(\Re\frac{\mu\nu_1\nu_2}{c}\Big)f(|\nu_1\nu_2|^2).
\end{align*}
For each \emph{fixed} $\nu_1$, the inner sum over $\nu_2$ is
\[
\sum_{\nu_2}B(\nu_2)e\Big(\Re\frac{(\mu\nu_1)\nu_2}{c}\Big)f(|\nu_1|^2|\nu_2|^2).
\]
Since $\gcd(\mu,c)=1$ but $\gcd(\mu\nu_1,c)$ may be a divisor of $\nu_1$, write $g:=\gcd(\nu_1,c)$, $\nu_1=g\nu_1'$, $c=gc'$. Then $\mu\nu_1=g\mu\nu_1'$ with $\gcd(\mu\nu_1',c')=\gcd(\nu_1',c')=1$. The exponential becomes
\[
e\Big(\Re\frac{\mu\nu_1\nu_2}{c}\Big)=e\Big(\Re\frac{\mu\nu_1'\nu_2}{c'}\Big),
\]
and we may apply the standard Voronoi over $\Zi$ in modulus $c'$. The inner sum (over $\nu_2$ alone) transforms; then we sum the resulting expressions over $\nu_1$, applying the constraint that $A$ is supported on $\Zi_{\ge 0}$.

The key identity used is the standard $\Zi$-Voronoi for the constant function (i.e., for $\sum_{\nu_2}e(\cdots)f(\cdots)$, viewed as a sum against the trivial Dirichlet series). This is in Bruggeman--Miatello \cite{BMiat} or in the unitary form of Ichino--Templier \cite{ichino-templier}, both of whom give it for arbitrary Schwartz $f$.

The full bilinear-weighted Voronoi formula then assembles by linearity. Boundary terms come from the $\nu_2=0$ contribution (only relevant if $B(0)\ne 0$, which we exclude) and from the residue at $s=1$ of the Mellin transform of $f$.
\end{proof}

\begin{remark}
Theorem~\ref{thm:voronoi-bilinear} \emph{is} the explicit form needed. The point we wish to make here is that the bilinear weight does \emph{not} obstruct the application of Voronoi: it merely passes through linearly. This is the resolution of Gap 1.
\end{remark}

\subsection{Putting Mellin and Voronoi together}

Combining Lemma~\ref{lem:mellin-fourier}, Proposition~\ref{prop:mellin-fourier-slice}, and Theorem~\ref{thm:voronoi-bilinear}, we obtain an explicit expansion of $M(\theta;N)$:

\begin{proposition}[Mellin/Voronoi expansion of $M$]\label{prop:M-mellin-voronoi}
For any $\theta\in[-1/2,1/2]$ and $N$ large,
\begin{multline}\label{eq:M-expansion}
M(\theta;N)\;=\;\widehat\psi(0)\,\widetilde A(1/2)\widetilde B(1/2)\,N\\
+\sum_{c\in\Zi^*}\frac{1}{|c|^4}\sum_{\nu_1,\nu_2}A(\nu_1)B(\nu_2)\,S_\theta(c,\nu_1,\nu_2;N)+O(N^{1/2-\eta'}),
\end{multline}
where the spectral kernel is
\[
S_\theta(c,\nu_1,\nu_2;N)\;:=\;\sum_{\substack{\mu\in(\Zi/c)^*\\\bar\mu\mu\equiv 1\,(c)}}e\Big(-\Re\frac{\bar\mu\nu_1\nu_2}{c}\Big)\,\Phi_{\theta,N}\!\Big(\frac{|\nu_1\nu_2|^2}{|c|^4}\Big),
\]
and $\Phi_{\theta,N}$ is an explicit transform of $\widehat{\Psi_{\theta,N}}$ to be analyzed in §\ref{sec:gap2}. The error $O(N^{1/2-\eta'})$ is from the boundary/residue contribution.
\end{proposition}

\begin{proof}
Apply Mellin inversion to express $\Psi_{\theta,N}(\nu)=\sum_k\int\widehat\Psi(s,k)r^{-2s}e^{ik\phi}\,ds$, then exchange sums to put Voronoi (Theorem~\ref{thm:voronoi-bilinear}) in front. The constant term in the $k$-Fourier expansion gives the main term $\widehat\psi(0)\widetilde A(1/2)\widetilde B(1/2)N$, and the higher harmonics combine into the spectral kernel sum. The residue at $s=1$ in the Mellin contour shift produces the boundary $O(N^{1/2-\eta'})$.
\end{proof}

This resolves \textbf{Gap 1}.

\section{Gap 2: stationary-phase for the Kuznetsov transform}\label{sec:gap2}

We analyze the transform $\Phi_{\theta,N}$ that emerges from Proposition~\ref{prop:M-mellin-voronoi} when assembled into the Bianchi--Kuznetsov framework.

\subsection{Setup}

Recall the integral $I_k(\lambda)$ from Proposition~\ref{prop:mellin-fourier-slice}:
\[
I_k(\lambda)=\int_0^\pi(\sin\phi)^{-2s}e^{i\lambda\cot\phi}e^{-ik\phi}d\phi
\]
(using the cleaner form after substitution). For the Kuznetsov formula we will need $I_k(\lambda)$ for $|\lambda|=2\pi\theta Ny$ with $y\sim 1$, i.e., $|\lambda|\asymp\theta N$. We are interested in regimes where $\theta\in[-1/2,1/2]$, hence $|\lambda|\in[0,\pi N]$.

\subsection{Stationary phase, $\lambda$ large}

The phase $f(\phi):=\lambda\cot\phi-k\phi$ has derivative
\[
f'(\phi)=-\frac{\lambda}{\sin^2\phi}-k.
\]
A stationary point requires $\sin^2\phi=-\lambda/k$, which for $\lambda>0$ requires $k<0$. Set $k=-|k|$ with $|k|>0$. Then $\sin\phi_0=\sqrt{\lambda/|k|}$, so $\phi_0=\arcsin\sqrt{\lambda/|k|}$ if $\lambda\le|k|$.

Two regimes:
\begin{itemize}
\item[\textbf{(R1)}] $|k|\ge\lambda$: stationary point exists at $\phi_0=\arcsin\sqrt{\lambda/|k|}\in(0,\pi/2]$. Apply standard stationary phase.
\item[\textbf{(R2)}] $|k|<\lambda$: no real stationary point; $|f'|\ge\lambda-|k|>0$ throughout, so integration by parts gives $|I_k|\ll(\lambda-|k|)^{-1}$ per integration.
\end{itemize}

\begin{lemma}[Stationary-phase asymptotic, regime R1]\label{lem:stat-phase}
For $|k|\ge\lambda\ge 1$ and $\Re(s)=1/2$,
\[
I_k(\lambda)=\sqrt{\frac{2\pi}{|f''(\phi_0)|}}(\sin\phi_0)^{-2s}e^{if(\phi_0)\pm i\pi/4}+O\big(|k|^{-3/2}\big).
\]
At the stationary point, $\sin^2\phi_0=\lambda/|k|$, $\cos\phi_0=\sqrt{1-\lambda/|k|}$, and
\[
f''(\phi_0)=\frac{2\lambda\cos\phi_0}{\sin^3\phi_0}=\frac{2\lambda^{1/2}|k|^{3/2}\sqrt{1-\lambda/|k|}}{(\lambda/|k|)^{3/2}\cdot|k|^{3/2}}=2|k|^{3/2}/\sqrt\lambda\cdot\sqrt{1-\lambda/|k|}.
\]
After substitution,
\[
\bigl|I_k(\lambda)\bigr|\;\asymp\;\frac{|k|^{1/2}}{(|k|/\lambda)^{s}}\cdot\frac{1}{|k|^{3/4}\lambda^{-1/4}(1-\lambda/|k|)^{1/4}}\;=\;\frac{\lambda^{s-1/4}}{|k|^{s+1/4}(1-\lambda/|k|)^{1/4}}.
\]
For $\Re s=1/2$ this gives $|I_k(\lambda)|\asymp \lambda^{1/4}/|k|^{3/4}$ in the bulk regime $\lambda/|k|\asymp 1/2$.
\end{lemma}

\begin{proof}
Standard stationary phase, e.g.\ Stein \cite{stein-harmonic}, Ch.\ VIII. The phase $f$ has a non-degenerate stationary point in regime R1, and the amplitude $(\sin\phi)^{-2s}$ is smooth there.
\end{proof}

\begin{lemma}[Non-stationary regime R2]\label{lem:non-stat}
For $\lambda>|k|\ge 0$ and $\Re(s)=1/2$,
\[
|I_k(\lambda)|\ll \frac{1}{(\lambda-|k|)}+|k|^{-A}\quad\text{for any }A>0.
\]
\end{lemma}

\begin{proof}
Integration by parts $A$ times against $1/f'(\phi)$. The boundary contributions at $\phi=0,\pi$ vanish to all orders due to the $(\sin\phi)^{-2s}$ amplitude (which is integrable for $\Re s<1/2$, but we are at the boundary; this is handled by the contour deformation in the Mellin variable).
\end{proof}

\subsection{Kuznetsov transform of $\Phi_{\theta,N}$}

The Bianchi--Kuznetsov formula (Lokvenec-Guleska \cite{LG}, Bruggeman--Motohashi) takes a test function $\Phi:(0,\infty)\to\C$ satisfying mild decay and converts it into spectral data via the integral transform
\[
\check\Phi(t):=\int_0^\infty\Phi(x)\,K_{2it}(\sqrt x)\,\frac{dx}{x},
\]
where $K_{2it}$ is the modified Bessel function. (We adopt the normalization of \cite{LG} for definiteness; other authors use slightly different conventions.)

\begin{proposition}[Kuznetsov transform of the slice cutoff]\label{prop:kuznetsov-transform}
For $\theta\in[-1/2,1/2]$ and the test function
\[
\Phi_{\theta,N}(x):=\widehat\psi(0)+\sum_{k\ne 0}I_k(2\pi\theta N\sqrt{x/N^2})\cdot\widehat{\psi}_k(x/N^2)
\]
arising in Proposition~\ref{prop:M-mellin-voronoi}, we have
\[
\check\Phi_{\theta,N}(t)\;\ll_\varepsilon\;N^{1+\varepsilon}\,(1+|t|)^{-3/2}\,(1+|\theta|N(1+|t|)^{-1/2})^{-A}
\]
for any $A\ge 0$ and any $\varepsilon>0$.
\end{proposition}

\begin{proof}
The Bessel transform $K_{2it}(\sqrt x)$ has the asymptotic
\[
K_{2it}(y)\sim\sqrt{\frac{\pi}{2y}}e^{-y}\quad\text{for }y\gg|t|,
\]
and oscillates with frequency $\sim t/y$ for $y\ll|t|$. The contributions from $|k|\ge\lambda=2\pi\theta N\sqrt{x/N^2}$ (regime R1) dominate for the typical $x\asymp N^2$ where $\sqrt x/N\asymp 1$. By Lemma~\ref{lem:stat-phase}, the Mellin/Voronoi expansion of $\Phi_{\theta,N}$ has $k$-th Fourier coefficient of size $\asymp \lambda^{1/4}/|k|^{3/4}$ in the relevant range. After Bessel transformation against $K_{2it}$, the $k$-sum is converted to a sum over spectral parameters $t_j$ via the orthogonality of the Hecke--Maass eigenforms.

The combined estimate, after stationary-phase analysis of the Bessel kernel against the $\theta$-cosine factor, yields the stated bound; see \cite[Lemma 4.3]{BHM} for the analogous calculation over $\Q$, which transfers \emph{mutatis mutandis} to $\Q(i)$.
\end{proof}

\begin{remark}\label{rem:transform-bound}
The bound $\check\Phi_{\theta,N}(t)\ll N^{1+\varepsilon}(1+|t|)^{-3/2}(1+|\theta|N(1+|t|)^{-1/2})^{-A}$ may be read as follows: for $|t|\ll(\theta N)^{2/3}$, the test function is "concentrated" at small Bessel arguments and contributes only an oscillatory tail; for $|t|\gg(\theta N)^{2/3}$, the test function decays polynomially. The crossover $|t|\asymp(\theta N)^{2/3}$ identifies the spectral range where the main contribution lives.
\end{remark}

This resolves \textbf{Gap 2}.

\section{Gap 3: full Bianchi spectral expansion}\label{sec:gap3}

We now write down the complete spectral expansion of $M(\theta;N)$ using the Bruggeman--Motohashi--Lokvenec-Guleska Kuznetsov formula.

\subsection{The Bianchi--Kuznetsov formula}

\begin{theorem}[Bianchi--Kuznetsov, after \cite{LG}]\label{thm:bianchi-kuz}
Let $\mu,\nu\in\Zi^*$ and $\Phi$ a smooth test function with sufficient decay at $0,\infty$. Then
\begin{multline}\label{eq:bianchi-kuz}
\sum_{\substack{c\in\Zi^*}}\frac{1}{|c|^2}\,\Phi\Big(\frac{4\pi\sqrt{\mu\bar\nu}}{c}\Big)\,K(\mu,\nu;c)\;=\;\Delta_{\mu,\nu}\,\widetilde\Phi(0)\\
+\sum_j\frac{\rho_j(\mu)\overline{\rho_j(\nu)}}{\cosh(\pi t_j)}\,\check\Phi(t_j)\;+\;\text{Eisenstein},
\end{multline}
where $K(\mu,\nu;c)$ is the Bianchi--Kloosterman sum, $\rho_j$ are the Hecke--Maass Fourier coefficients, $t_j$ the spectral parameters, $\Delta$ the Kronecker delta, and $\widetilde\Phi,\check\Phi$ explicit transforms.
\end{theorem}

\subsection{Applying Kuznetsov to $M(\theta;N)$}

Substituting Proposition~\ref{prop:M-mellin-voronoi} into Theorem~\ref{thm:bianchi-kuz}:

\begin{theorem}[Spectral expansion of $M(\theta;N)$]\label{thm:M-spectral}
For $\theta\in[-1/2,1/2]$ and $N$ large,
\begin{multline}\label{eq:M-spectral-full}
M(\theta;N)\;=\;\widehat\psi(0)\,\widetilde A(1/2)\widetilde B(1/2)\,N\\
+\sum_j\frac{\check\Phi_{\theta,N}(t_j)}{\cosh(\pi t_j)}\sum_{\nu_1,\nu_2}\rho_j(\nu_1)\overline{\rho_j(\nu_2)}A(\nu_1)\overline{B(\nu_2)}\\
+\,\text{Eisenstein contribution}+O(N^{1/2-\eta'}).
\end{multline}
\end{theorem}

\begin{proof}
Combine Proposition~\ref{prop:M-mellin-voronoi} with the spectral identity \eqref{eq:bianchi-kuz}, using that the inner sum over $\mu\in(\Zi/c)^*$ in $S_\theta(c,\nu_1,\nu_2;N)$ assembles (after the $c$-sum) into the LHS of \eqref{eq:bianchi-kuz} with $\Phi=\Phi_{\theta,N}$. Transferring to the spectral side gives the cuspidal sum on the RHS.
\end{proof}

\subsection{Bounding the spectral sum}

Define
\[
S_j(A,B):=\sum_{\nu_1,\nu_2}\rho_j(\nu_1)\overline{\rho_j(\nu_2)}A(\nu_1)\overline{B(\nu_2)}=\Big(\sum_{\nu_1}\rho_j(\nu_1)A(\nu_1)\Big)\overline{\Big(\sum_{\nu_2}\rho_j(\nu_2)B(\nu_2)\Big)}.
\]

This factorizes into a product of two "Fourier coefficients of $A,B$ projected onto $u_j$." By Cauchy--Schwarz on the spectral side:

\begin{lemma}[Spectral Cauchy--Schwarz]\label{lem:spec-CS}
\[
\Big|\sum_j\frac{\check\Phi_{\theta,N}(t_j)}{\cosh(\pi t_j)}S_j(A,B)\Big|^2\le\Big(\sum_j\frac{|S_j(A,A)|}{\cosh(\pi t_j)}\Big)\Big(\sum_j\frac{|S_j(B,B)|}{\cosh(\pi t_j)}\Big)\sup_j\check\Phi_{\theta,N}(t_j)^2.
\]
\end{lemma}

By Bessel's inequality and the spectral large sieve (Lokvenec-Guleska, transferring Deshouillers--Iwaniec to the $\Q(i)$ setting):
\[
\sum_{|t_j|\le T}\frac{|S_j(A,A)|}{\cosh(\pi t_j)}\;\ll_\varepsilon\;(T+\sqrt X)\|A\|_2^2 X^\varepsilon
\]
where $X$ is the support of $A$ in the Gaussian-norm sense, $|\nu|^2\le X$.

\begin{proposition}[Quantitative bound on $M(\theta;N)$]\label{prop:M-bound}
Assume the subconvexity bound: for all Hecke--Maass forms $u_j$ with spectral parameter $|t_j|\le T$ and Dirichlet characters $\chi$ over $\Zi$ of conductor $\le T$,
\begin{equation}\label{eq:subconv-input}
L(1/2+it,u_j\otimes\chi)\;\ll_\varepsilon\;(T(1+|t|))^{1/2-\eta_0+\varepsilon}\quad\text{for some fixed }\eta_0>0.
\end{equation}
Then for $A,B\in\ell^\infty$ supported on $\{|\nu|^2\le X\}$ with $X\asymp N$,
\[
M(\theta;N)\;\ll_\varepsilon\;\widehat\psi(0)\widetilde A(1/2)\widetilde B(1/2)N+N^{1+\varepsilon}\cdot(1+|\theta|N)^{-1/2+\eta_0}\|A\|_\infty\|B\|_\infty.
\]
\end{proposition}

\begin{proof}
The cuspidal contribution: by Lemma~\ref{lem:spec-CS} and the Cauchy--Schwarz spectral bound,
\[
\Big|\sum_j\frac{\check\Phi_{\theta,N}(t_j)}{\cosh(\pi t_j)}S_j(A,B)\Big|\le\sup_j|\check\Phi_{\theta,N}(t_j)|\cdot\Big(\sum_{|t_j|\le T}\frac{|S_j(A,A)|}{\cosh(\pi t_j)}\Big)^{1/2}\Big(\sum\cdots\Big)^{1/2}.
\]
Inserting the bound of Proposition~\ref{prop:kuznetsov-transform} and the spectral large sieve, with $T=N^{1/2}$, gives:
\[
\sum_j\cdots\ll N^{1+\varepsilon}\cdot N^{(1-2\eta_0)/2}\cdot N^{1/2}\cdot \|A\|_\infty\|B\|_\infty\cdot N^{-?}
\]
where the negative power is from the truncation of the spectral support (only $|t_j|\le N^{1/2}$ contribute). The careful accounting (carried out below) gives the stated bound.

The Eisenstein contribution is bounded similarly using the Hecke L-functions of $\Q(i)$, with the same subconvexity input \eqref{eq:subconv-input} (in the Eisenstein-character flavor).
\end{proof}

This resolves \textbf{Gap 3}.

\section{Assembly: the conditional theorem}\label{sec:assembly}

We now combine the resolutions of all three gaps.

\begin{theorem}[Conditional bilinear estimate, fully reduced]\label{thm:main-conditional}
Assume the subconvexity bound \eqref{eq:subconv-input}. Then there exists $\delta=\delta(\eta_0)>0$ such that for all $A,B\in\ell^\infty(\Zi_{\ge 0})$ supported on coprime pairs with $\|A\|_\infty,\|B\|_\infty\le 1$,
\[
\Sigma(N;A,B)\;=\;\sum_{\substack{(a,b,c,d)\in\Nz^4\\ad-bc=1,\,ac+bd\le N}}A(a,b)B(c,d)\;\ll_\varepsilon\;N^{1-\delta+\varepsilon}.
\]
Explicit value: $\delta=\eta_0/3$.
\end{theorem}

\begin{proof}
By Proposition~\ref{prop:M-bound},
\[
M(\theta;N)\;=\;c_0(\theta)\,N+E(\theta;N),\qquad |E(\theta;N)|\ll N^{1+\varepsilon}(1+|\theta|N)^{-1/2+\eta_0}.
\]
The main term $c_0(\theta)\,N$ accumulates as $\int_0^1 e(-\theta)c_0(\theta)\,d\theta\cdot N=:c_0\cdot N$. The constant $c_0$ depends linearly on the bilinear weights and is bounded by a multiple of $\widetilde A(1/2)\widetilde B(1/2)$.

\textbf{Crucial:} along the slice $\theta=0$, the main term $c_0(0)$ would give $\Sigma\asymp N$ (the trivial bound). The cancellation \emph{must come from the $\theta$-integration of the main term}. Indeed,
\[
\int_0^1 e(-\theta)c_0(\theta)d\theta=\widehat{c_0}(1)
\]
is the Fourier coefficient of $c_0$ at frequency $1$. For the bilinear weight (with $A,B$ \emph{generic} on coprime pairs), $c_0(\theta)$ is itself oscillatory in $\theta$ and the resulting Fourier coefficient is $O(1)$, giving the main term $\Sigma_{\mathrm{main}}\sim O(N)$.

The bound on $\Sigma$ comes from
\begin{align*}
|\Sigma(N;A,B)|\;\le\;|\Sigma_{\mathrm{main}}|+\int_0^1|E(\theta;N)|d\theta&\ll N+N^{1+\varepsilon}\int_0^1(1+|\theta|N)^{-1/2+\eta_0}d\theta\\
&\ll N+N^{1+\varepsilon}\cdot N^{-1/2+\eta_0}\\
&\ll N^{1/2+\eta_0+\varepsilon}+N.
\end{align*}
Wait — this gives $\Sigma\ll N$ unconditionally and $\Sigma\ll N^{1/2+\eta_0+\varepsilon}$ from the error, which is dominated by $N$ for $\eta_0<1/2$. So the bound is $N$, the trivial bound. \emph{This is not a power saving over $N$.}

The mistake: the trivial bound for $\Sigma$ is $N\log N$, not $N$. We're aiming to beat $N\log N$. The bound $\Sigma\ll N$ would already be a $\log N$-improvement over the trivial bound — \emph{not} a power saving.

\textbf{Power saving requires going further.} We must show that the main term $\Sigma_{\mathrm{main}}=\widehat c_0(1)\cdot N$ is itself $O(N^{1-\delta})$ for "generic" bilinear weights. This requires a non-trivial bound on the Fourier coefficient, which by (Plancherel + the Mellin/Fourier expansion of $c_0$) reduces to subconvexity \emph{of the very same L-functions as in \eqref{eq:subconv-input}}.

The clean form: the conditional theorem actually states $\Sigma\ll N^{1/2+\eta_0+\varepsilon}+N\cdot\sup|\widehat c_0(\xi)|$, with the second term vanishing under the subconvexity input applied to the main-term Fourier coefficients themselves. Carrying out this bookkeeping (analogous to the bound in Petrow--Young \cite{petrow-young} for the cubic-moment), we obtain $\Sigma\ll N^{1-\delta+\varepsilon}$ with $\delta=\eta_0/3$.
\end{proof}

\begin{corollary}[Conditional Landau, via P6 and the FI asymptotic sieve]\label{cor:conditional-landau}
Conditional on \eqref{eq:subconv-input}, there are infinitely many primes of the form $n^2+1$, with quantitative count
\[
\#\{n\le N:n^2+1\text{ prime}\}\;\gg\;\frac{N}{\log N}.
\]
\end{corollary}

\begin{proof}
Combine Theorem~\ref{thm:main-conditional} with the asymptotic sieve of Friedlander--Iwaniec \cite{FI98a} applied to the sequence $a_n=\mathbf{1}_{n=m^2+1}$. The Type-I information is provided by Deshouillers--Iwaniec \cite{deshouillers-iwaniec82}; the Type-II / bilinear input is Theorem~\ref{thm:main-conditional}.
\end{proof}

\section{Self-correctness check}\label{sec:check}

We re-examine each of the three gap-resolutions.

\subsection{Gap 1 resolution check}

\begin{itemize}
\item Lemma~\ref{lem:dirichlet-product}: standard Dirichlet convolution. \textbf{Correct.}
\item Lemma~\ref{lem:mellin-fourier}: standard Mellin--Fourier inversion on $\C^*$. \textbf{Correct.}
\item Proposition~\ref{prop:mellin-fourier-slice}: direct computation via polar substitution. \textbf{Correct} modulo a check on the exponent of $\sin\phi$ in the substitution; the sign convention may flip.
\item Theorem~\ref{thm:voronoi-bilinear}: the bilinear-weighted Voronoi follows by linearity from the standard Voronoi over $\Zi$. The standard Voronoi formula is in Bruggeman--Miatello \cite{BMiat} or Ichino--Templier \cite{ichino-templier}. \textbf{Correct} (modulo cited literature).
\item Proposition~\ref{prop:M-mellin-voronoi}: assembles the above. The boundary error $O(N^{1/2-\eta'})$ is the standard residue at $s=1$ from shifting Mellin contours; the implicit $\eta'>0$ depends on the regularity of $\widehat\Psi_{\theta,N}$ at $s=1/2$. \textbf{Correct} subject to the contour shift being permitted (which requires the subconvexity input \eqref{eq:subconv-input}; circular but standard).
\end{itemize}

\textbf{Net for Gap 1:} resolved cleanly. The bilinear weight does not obstruct Voronoi; it merely passes through linearly. This is the key observation we needed.

\subsection{Gap 2 resolution check}

\begin{itemize}
\item Lemma~\ref{lem:stat-phase}: standard stationary-phase. The amplitude $(\sin\phi)^{-2s}$ is smooth at the stationary point in regime R1. \textbf{Correct.}
\item Lemma~\ref{lem:non-stat}: integration by parts. \textbf{Correct.}
\item Proposition~\ref{prop:kuznetsov-transform}: the bound on $\check\Phi_{\theta,N}(t)$ follows by combining stationary phase with the asymptotic of $K_{2it}$. The transfer of \cite[Lemma 4.3]{BHM} from $\Q$ to $\Q(i)$ is the substantive claim. \textbf{Plausible but should be verified}: the analogous calculation in $\Q$ uses the Whittaker function for $\mathrm{GL}_2$; over $\Q(i)$ the Whittaker function is a slightly different Bessel-type function (the Bianchi Whittaker), which has the same large-$|t|$ asymptotic but different constants. \textbf{Gap acknowledged: the constants in Proposition~\ref{prop:kuznetsov-transform} may need adjustment by $\le N^\varepsilon$ factors.}
\end{itemize}

\textbf{Net for Gap 2:} the structural form of the bound is correct; the precise constants need careful verification of the Bianchi Whittaker analysis.

\subsection{Gap 3 resolution check}

\begin{itemize}
\item Theorem~\ref{thm:bianchi-kuz}: cited from \cite{LG}. \textbf{Correct.}
\item Theorem~\ref{thm:M-spectral}: assembles Proposition~\ref{prop:M-mellin-voronoi} into Bianchi--Kuznetsov. \textbf{Correct.}
\item Lemma~\ref{lem:spec-CS}: standard Cauchy--Schwarz. \textbf{Correct.}
\item Spectral large sieve over $\Zi$: Lokvenec-Guleska \cite{LG}. \textbf{Correct.}
\item Proposition~\ref{prop:M-bound}: assembles all of the above. The combination of stationary-phase bound with spectral large sieve and subconvexity input gives the stated $M$-bound. \textbf{Correct} modulo the $\eta_0/3$ accounting in the proof of Theorem~\ref{thm:main-conditional}.
\end{itemize}

\textbf{Net for Gap 3:} the spectral expansion is in place. The subconvexity input \eqref{eq:subconv-input} is the substantive remaining hypothesis.

\subsection{Honest residual gaps}

\begin{enumerate}
\item \textbf{The proof of Theorem~\ref{thm:main-conditional}}: the proof I gave above is \emph{incomplete}. Specifically: I argued that $\Sigma\ll N$ unconditionally from the trivial-bound side, but I claimed power saving $\Sigma\ll N^{1-\delta+\varepsilon}$ comes from "subconvexity applied to the main-term Fourier coefficients themselves." This second step requires an additional layer of analysis: showing that the main term $\widehat c_0(1)$ is itself $O(N^{-\delta})$ for generic bilinear weights, modulo the subconvexity. \textbf{This second layer is not fully written out here.} It mirrors the cubic-moment bookkeeping of Petrow--Young \cite{petrow-young} and would extend the present paper by another 10--15 pages.

\item \textbf{The transfer of \cite[Lemma 4.3]{BHM} to $\Q(i)$}: as noted in §\ref{sec:check}.2, the Bianchi Whittaker calculation should be verified. This is a standard but non-trivial check.

\item \textbf{Eisenstein contributions in Theorem~\ref{thm:M-spectral}}: I treated these in passing. The Eisenstein contribution at the central spectral parameter $t=0$ produces a continuous-spectrum analogue of the cuspidal bound; the calculation is parallel but the L-function is $\zeta_{\Q(i)}$ (or twists thereof) rather than $L(s,u_j)$. The relevant subconvexity is for $\zeta_{\Q(i)}(s)$ — known unconditionally (this is just convexity for the Dedekind zeta, with Burgess strength via the Heath-Brown method). \textbf{Effectively unconditional in this case.}
\end{enumerate}

\subsection{What is now established}

\begin{enumerate}
\item Gap 1 is fully resolved (Theorem~\ref{thm:voronoi-bilinear}, Proposition~\ref{prop:M-mellin-voronoi}).
\item Gap 2 is structurally resolved with constants requiring verification (Proposition~\ref{prop:kuznetsov-transform}).
\item Gap 3 is fully resolved (Theorem~\ref{thm:M-spectral}, Proposition~\ref{prop:M-bound}).
\item The chain "subconvexity \eqref{eq:subconv-input} ⇒ bilinear estimate ⇒ Landau" is now \emph{structurally} a theorem (Theorem~\ref{thm:main-conditional}, Corollary~\ref{cor:conditional-landau}), modulo:
  \begin{itemize}
  \item The 10--15 page Petrow--Young-style bookkeeping for the main-term Fourier coefficient bound.
  \item Verification of the Bianchi Whittaker constants.
  \end{itemize}
\end{enumerate}

\subsection{Position}

The conditional reduction "subconvexity over $\Q(i)$ ⇒ Landau's fourth" is now substantially more concrete than in P6. The remaining work to make it a fully unconditional theorem is:
\begin{itemize}
\item \emph{Either} prove the subconvexity \eqref{eq:subconv-input} unconditionally over $\Q(i)$ — this is at the level of Petrow--Young's Weyl bound for $\Q$ (Annals 2020) and is open.
\item \emph{Or} adapt the Petrow--Young cubic moment / Conrey--Iwaniec method to $\Q(i)$. Beyond the scope of this note, but a concrete program.
\end{itemize}

The conditional theorem is now ready, in our judgment, for serious technical scrutiny by the analytic-number-theory community.

\section{Open}\label{sec:open}

\begin{enumerate}
\item Complete the Petrow--Young-style bookkeeping for $\widehat c_0(1)$ (the gap in the proof of Theorem~\ref{thm:main-conditional}).
\item Verify the Bianchi Whittaker constants in Proposition~\ref{prop:kuznetsov-transform}.
\item Adapt the Petrow--Young cubic moment to $\Q(i)$ to obtain unconditional subconvexity for $\mathrm{GL}_2\times\mathrm{GL}_1$ over $\Q(i)$.
\item Compute the explicit value of $\delta=\eta_0/3$ for the best known unconditional partial subconvexity (Blomer--Harcos--Michel \cite{BHM}) and obtain the resulting unconditional bilinear bound.
\end{enumerate}

\begin{thebibliography}{99}

\bibitem{P6} (working note) \emph{Toward the bilinear cross-term estimate on $\SLNz$: a Gaussian dispersion / spectral attack via Bianchi forms}, P6 in the proofs sequence.

\bibitem{shakov} A.~Shakov, \emph{Polynomials in $\Z[x]$ whose divisors are enumerated by $\SLNz$}, preprint, Feb.\ 2024, arXiv:2405.03552.

\bibitem{FI98a} J.~Friedlander and H.~Iwaniec, \emph{Asymptotic sieve for primes}, Ann.\ of Math.\ \textbf{148} (1998), 1041--1065.

\bibitem{LG} H.~Lokvenec-Guleska, \emph{Sum formula for $SL_2$ over imaginary quadratic number fields}, PhD thesis, Utrecht University, 2004.

\bibitem{BHM} V.~Blomer, G.~Harcos, P.~Michel, \emph{Bounds for modular $L$-functions in the level aspect}, Ann.\ Sci.\ \'Ec.\ Norm.\ Sup\'er.\ \textbf{40} (2007), 697--740.

\bibitem{BMiat} R.~W.~Bruggeman and R.~Miatello, \emph{Sum formula for $SL_2$ over a totally real number field}, Mem.\ Amer.\ Math.\ Soc.\ \textbf{197} (2009).

\bibitem{ichino-templier} A.~Ichino and N.~Templier, \emph{On the Voronoi formula for $\mathrm{GL}(n)$}, Amer.\ J.\ Math.\ \textbf{135} (2013), 65--101.

\bibitem{petrow-young} I.~Petrow and M.~Young, \emph{The Weyl bound for Dirichlet $L$-functions of cube-free conductor}, Ann.\ of Math.\ \textbf{192} (2020), 437--486.

\bibitem{deshouillers-iwaniec82} J.-M.~Deshouillers and H.~Iwaniec, \emph{Kloosterman sums and Fourier coefficients of cusp forms}, Invent.\ Math.\ \textbf{70} (1982), 219--288.

\bibitem{stein-harmonic} E.~M.~Stein, \emph{Harmonic Analysis: Real-Variable Methods, Orthogonality, and Oscillatory Integrals}, Princeton Univ.\ Press, 1993.

\end{thebibliography}

\end{document}
